\subsection{传统功法的安全审视：铁裆功、气功导引与提肛}

中医性养生中的功法训练近年随短视频传播走红，其中一些环节需要区分"有证据的部分"与"有风险的部分"：

\begin{itemize}
\item \textbf{提肛（撮谷道）}：主动收缩盆底肌的传统练法，与西医凯格尔训练（Kegel）同源——这是本书明确支持的部分，对射精控制、勃起硬度与女性盆底健康都有循证支持（参见盆底肌训练相关章节）。"日撮百遍"的频次表述与凯格尔方案的常规剂量也相容，中西医结合，殊途同归；
\item \textbf{气功导引与呼吸法}：房中文献记载的交合呼吸法（深吸缓呼、入静放松、以意导气），其实质是副交感神经激活训练——与行为治疗中的放松训练、延缓射精的呼吸调节相通。作为放松技巧安全且可取；需要剥离的是其后附加的"还精补脑""采阴补阳"等玄学承诺；
\item \textbf{铁裆功}：传统练法包含对睾丸、会阴部的揉捻、吊挂与负重，宣称"固精强肾"。需要明确：重物坠吊与暴力揉捏没有现代证据支持，反而存在睾丸挫伤、睾丸扭转与会阴神经损伤的风险——\textbf{睾丸扭转是急症}，拖延处理可致睾丸坏死（参见男性卷急症章节）。对这类内容，"传统文化"不是安全认证；
\item \textbf{警惕商业包装}：以功法为名的天价"壮阳课程""私教补肾班"是养生骗局的常见载体，话术往往混合真实的盆底训练与伪科学承诺——凡是承诺"根治""增长""再发育"的课程，都值得先复习本书勃起生理与阴茎解剖章节，再决定要不要交钱。
\end{itemize}
